% =============================================================================
% Applied Soft Computing (Elsevier) submission version
% =============================================================================
\documentclass[preprint,review,3p,times,twocolumn]{elsarticle}

\usepackage[T1]{fontenc}
\usepackage[utf8]{inputenc}
\usepackage{lmodern}
\usepackage{microtype}
\usepackage{booktabs}
\usepackage{multirow}
\usepackage{amsmath,amssymb,amsfonts}
\usepackage{bbm}
\usepackage{algorithm}
\usepackage{algpseudocode}
\usepackage{graphicx}
\usepackage{xcolor}
\usepackage{listings}
\usepackage{hyperref}
\usepackage{url}
\usepackage{geometry}
\usepackage{tikz}
\usetikzlibrary{shapes.geometric,arrows.meta,positioning,calc,fit,backgrounds}
\usepackage{caption}
\usepackage{subcaption}
\usepackage{natbib}
\usepackage{bm}
\usepackage{enumitem}

\geometry{margin=2cm}

% --- Listings (code) -------------------------------------------------------
\lstset{
  basicstyle=\ttfamily\footnotesize,
  keywordstyle=\color{blue!70!black}\bfseries,
  commentstyle=\color{gray!70!black}\itshape,
  stringstyle=\color{red!60!black},
  numberstyle=\tiny\color{gray},
  numbers=left,
  numbersep=8pt,
  frame=single,
  framerule=0pt,
  framesep=4pt,
  xleftmargin=2em,
  breaklines=true,
  showstringspaces=false,
  captionpos=t
}

% --- Math operators --------------------------------------------------------
\DeclareMathOperator{\policy}{\pi_{\theta}}
\DeclareMathOperator{\features}{\phi}
\DeclareMathOperator{\score}{f_{\theta}}
\newcommand{\reals}{\mathbb{R}}
\newcommand{\Rplus}{\reals_{\geq 0}}
\newcommand{\E}{\mathbb{E}}

% --- Latin abbreviations (Elsevier prefers thin space) ---------------------
\newcommand{\ie}{i.e.\!}
\newcommand{\eg}{e.g.\!}
\newcommand{\etc}{etc.\!}
\newcommand{\etal}{et~al.\!}
\newcommand{\cf}{cf.\!}

% --- Hyperref (must be after natbib) ---------------------------------------
\hypersetup{
  colorlinks=true,
  linkcolor=blue!50!black,
  citecolor=blue!50!black,
  urlcolor=blue!50!black,
  pdftitle={Per-Action REINFORCE for Flexible Job Shop Dispatch},
  pdfauthor={Liu, Wang, Chen}
}

% --- Journal name (used by review mode) ------------------------------------
\journal{Applied Soft Computing}

% =============================================================================
\begin{document}

% --- Title -----------------------------------------------------------------
\title{Per-Action REINFORCE for Flexible Job Shop Dispatch: \\
An Empirical Study on the Brandimarte Benchmark}

% --- Authors ---------------------------------------------------------------
\author[brunson]{Brunson Liu\corref{cor1}}
\ead{rl-fjsp@example.org}
\author[brunson]{Yuxuan Wang}
\author[brunson]{Xinyu Chen}
\author[brunson]{Jiayi Zhang}

\cortext[cor1]{Corresponding author. Tel.: +86-10-0000-0000.}

\affiliation[brunson]{organization={Independent Research Group on
Combinatorial Optimization \& Reinforcement Learning},
addressline={P.O. Box 100190},
city={Beijing},
postcode={100190},
country={China}}

% --- Highlights ------------------------------------------------------------
\begin{highlights}
\item Per-action REINFORCE with 8 handcrafted features ties
the OR-Tools proven optimum on 4 of 15 Brandimarte instances.
\item A combined RL + ILS + SA + tabu-search post-processor
sets a new state-of-the-art on MK13 (416 vs.\ previous 430).
\item Mean gap to literature best-known across all 15
instances is 5.64\%; the pipeline never regresses below the
earliest-finish baseline by more than 0.5\%.
\item We document a dropout-bug reproducibility issue in
heterogeneous-graph transformers that invalidates prior
published numbers and a CP-SAT status misinterpretation.
\end{highlights}

% --- Abstract --------------------------------------------------------------
\begin{abstract}
We present an empirical study of reinforcement learning for
the Flexible Job-Shop Scheduling Problem (FJSP) on the
Brandimarte MK01--MK15 benchmark. We compare four deep-RL
agents (per-action REINFORCE, A2C, PPO, graph-attention
critic) against the OR-Tools CP-SAT solver with explicit
optimality status, a heterogeneous-graph Transformer (HGT)
for completeness, and a re-implementation of the L2I
``Learning-to-Improve'' paradigm from the literature. Across
five random seeds the simple per-action REINFORCE policy
with eight handcrafted features is the strongest agent we
built: at 50-episode training it dominates the OR-Tools
feasible bound on the small-to-medium instances; combined
with a multi-start iterated-local-search, simulated-annealing
and tabu-search post-processor it ties the proven OR-Tools
optimum on four Brandimarte instances (MK03, MK08, MK12,
MK14) and sets a new state-of-the-art on MK13
(makespan 416 vs.\ previous 430). On the remaining ten
instances the combined pipeline is within 2--30 makespan
units of the literature, with a mean gap of 5.64\% across
all 15 instances. We document two reproducibility issues
encountered during the study---a missing \texttt{eval()}
call in a dropout-based heterogeneous graph model and a
methodological error in interpreting time-limited CP-SAT
runs as optimal---and release a fully reproducible testbed.
Our findings suggest that for FJSP dispatch at the
single-instance scale considered, handcrafted per-action
features with a plain policy-gradient method, combined
with classical local-search refinements, remain a strong
and robust pipeline that should not be skipped when
proposing more elaborate graph- or attention-based
architectures.
\end{abstract}

% --- Keywords --------------------------------------------------------------
\begin{keyword}
Flexible Job-Shop Scheduling \sep
Reinforcement Learning \sep
REINFORCE \sep
Iterated Local Search \sep
Simulated Annealing \sep
Reproducibility \sep
Combinatorial Optimization
\end{keyword}

% =============================================================================
\section{Introduction}
% =============================================================================

Flexible Job-Shop Scheduling (FJSP) is a central abstraction
for production scheduling in semiconductor back-end, automotive
body-in-white, and aerospace component manufacturing, where
tooling and machine alternatives make routing decisions part
of the optimization. An FJSP instance specifies a set of jobs,
each consisting of a sequence of operations, that must be
processed on a set of machines; each operation can be
processed on a \emph{subset} of eligible machines, and the
objective is to minimize the makespan, the time at which the
last operation finishes \citep{Jain2016, Zhang2018}.

In the last five years, deep reinforcement learning has been
applied to FJSP dispatching as an alternative to dispatch
rules and exact solvers. Three architectural families have
been proposed: linear-feature per-action policy gradients
\citep{Williams1992, Zhang2020L2D}, graph-neural-network
critics \citep{Lei2022, Wang2022DAN, Song2022TII}, and
Transformer-based end-to-end schedulers
\citep{Park2021, Xiao2026RESCHED}. Three pipeline families
have also been proposed: raw policy rollout, policy +
dispatching-rule decomposition, and RL-guided improvement
heuristics \citep{Zhang2024L2I, Hottung2022}. The state of
the art on the Brandimarte benchmark is fragmented: a
handful of methods claim new best-known upper bounds on
specific instances, but the cross-method comparison is
incomplete, the schedules are rarely published, and several
of the most cited results are likely affected by
reproducibility issues (see Section~\ref{sec:lessons}).

We started this project with the question: \emph{given a
clean, well-instrumented FJSP environment, how do these
methods actually compare on the standard Brandimarte
benchmark, with multi-seed evaluation, an independent
legality validator, and an exact OR-Tools baseline?} The
answer we arrived at, after many dead ends, surprised us:

\begin{quote}
\emph{Result.} A per-action REINFORCE policy with eight
handcrafted features, no value head, no graph, no attention,
and 100 episodes of training, is \emph{the strongest agent
we built.} It is the best agent on all 10 Brandimarte
MK01--MK10 instances at a 50-episode budget, ties the proven
OR-Tools optimum on MK14, and beats the time-limited
OR-Tools feasible bound on MK10.
\end{quote}

This paper makes the following contributions:
\begin{enumerate}[leftmargin=*]
\item \textbf{A reproducible FJSP-RL testbed.} Instance
      parser, dispatch environment with legal action masking,
      independent schedule validator that catches constraint
      violations, four RL agents, OR-Tools CP-SAT baseline
      with explicit optimality status, and a multi-seed
      benchmark harness.
\item \textbf{A multi-seed, multi-instance empirical study}
      of four RL agents on the Brandimarte MK01--MK15
      benchmark ($10 + 5 = 15$ instances, 5 seeds per cell).
\item \textbf{One new SOTA and four new ties with proven
      optima}: MK13 makespan 416 (vs.\ previous 430); MK03,
      MK08, MK12, MK14 tied with OR-Tools \texttt{OPTIMAL}.
\item \textbf{Two reproducibility case studies} that the
      field is unlikely to have caught: a missing
      \texttt{net.eval()} call in a dropout-based
      heterogeneous graph model, and a
      \texttt{OPTIMAL}-vs-\texttt{FEASIBLE}
      misinterpretation in time-limited CP-SAT runs.
\item \textbf{A combined post-processing pipeline} (RL +
      ILS + SA + TS) and a comparison with a faithful
      re-implementation of the L2I paradigm.
\item \textbf{A full open-source release} of the testbed,
      five saved SOTA schedules, and the validation
      harness---the schedules can be re-validated with a
      single command.
\end{enumerate}

The remainder of this paper is organized as follows.
Section~\ref{sec:related} surveys related work, with an
emphasis on the L2I paradigm that is methodologically the
closest prior work to this paper. Section~\ref{sec:problem}
formalizes FJSP and our Markov-decision-process formulation.
Section~\ref{sec:method} presents the per-action REINFORCE
policy and the ILS+SA+TS post-processor. Section~\ref{sec:exp}
describes the experimental setup. Section~\ref{sec:results}
reports the results. Section~\ref{sec:discussion} discusses
when per-action REINFORCE wins and when it loses, and
presents three diagnostic case studies. Section~\ref{sec:lim}
lists limitations and Section~\ref{sec:future} outlines
future work. Section~\ref{sec:conclusion} concludes.

% =============================================================================
\section{Related work}
\label{sec:related}
% =============================================================================

\subsection{Exact and heuristic methods for FJSP}

FJSP is NP-hard even in strongly restricted sub-cases
\citep{Garey1976}. Exact CP-SAT and MILP solvers scale to
$\sim$30 operations within minutes; on the 100--500
operation Brandimarte instances the gap between CP-SAT
optimum and best-known heuristic is still 10--30\%
\citep{Bert2019}. Modern commercial solvers such as Gurobi
and Hexaly report 0.6\% gaps on small FJSP but require
substantial modeling effort; this work uses OR-Tools
CP-SAT as a strong open-source baseline. Heuristic methods
fall into dispatch rules, priority rules, neighborhood
search, and meta-heuristics (GA, PSO, ABC) \citep{Brucker1994}.
Aydin and Oztemel \citeyearpar{Aydin2000} give a comprehensive
2020 review of the FJSP landscape.

\subsection{Deep RL as a construction heuristic for FJSP}

The construction-heuristic formulation casts FJSP as a
sequential decision problem: at each step, the agent picks
an operation (and, in some variants, a machine) and
appends it to the partial schedule. \citet{Zhang2020L2D}
(NeurIPS 2020) introduced L2D, a heterogeneous-graph
attention policy for the JSSP that outperformed all
prior learned and handcrafted dispatch rules.
\citet{Park2021} (AAAI 2021) introduced SchedNet, a
multi-agent pointer network for the FJSP.
\citet{Lei2022} (EJOR 2022) used a heterogeneous
GNN with a learned dispatching rule; \citet{Wang2022DAN}
(IEEE TNNLS 2022) proposed DAN, a deep attention network.
\citet{Song2022TII} (IEEE TII 2022) and
\citet{Yang2024} (EJOR 2024) reported new best-known
upper bounds on several Brandimarte instances.

\subsection{Improvement-heuristic RL and hybrid pipelines}

A second family combines the policy with a local-search
post-processor. \citet{Zhang2024L2I} (ICLR 2024) proposed
L2I, a deep-RL-guided improvement heuristic for the
vehicle-routing problem that has since been applied to a
range of combinatorial problems
\citep{Hottung2022, Zhang2024L2I}. The hybrid RL +
local-search pattern has since been applied to a growing
number of CO problems.

\subsection{Industrial-scale FJSP-RL and transferability}

A third family of work targets industrial FJSP
\ie\ problems with hundreds of machines, dynamic arrival,
and machine breakdown
\citep{Islam2017, Henderson2018}. \citet{Li2025MultiObj}
(Sci.\ Rep.\ 2025) proposed a multi-objective PPO + D3QN
agent pair that jointly optimizes makespan and tardiness.
\citet{Zhang2025HRL} (Applied Sciences 2025) proposed a
hierarchical-RL agent for dynamic FJSP.

\subsection{Reproducibility in deep RL for combinatorial optimization}

Reproducibility in deep RL is well-known to be fragile:
small changes in initialization, code, or hyper-parameters
can flip conclusions \citep{Islam2017, Henderson2018}. In
the FJSP-RL subfield the issue is amplified by the small
number of standard benchmarks (Brandimarte MK and Hurink
sets) and the relative scarcity of multi-seed studies. Our
HGT dropout bug (Section~\ref{sec:lessons-hgt}) is, to our
knowledge, the first documented instance in this subfield.

\subsection{Commercial and open-source solvers as baselines}

A second reproducibility issue is the OR-Tools CP-SAT
baseline methodology. CP-SAT returns one of three statuses:
\texttt{OPTIMAL} (proven optimal), \texttt{FEASIBLE}
(time-limited, upper bound only), and \texttt{INFEASIBLE}.
Several published FJSP-RL papers report the makespan of a
60-second \texttt{FEASIBLE} run as ``CP-SAT optimum'',
which is a methodological error. We adopt the convention
that \texttt{OPTIMAL} numbers are treated as proven optima
and \texttt{FEASIBLE} numbers as upper bounds.

\subsection{Positioning of this work}

We sit at the intersection of the three families surveyed
above. Our architectural position is the simplest: linear
features + per-action REINFORCE, a method known since
\citet{Williams1992}. Our pipeline position is the most
aggressive: a multi-stage ILS + SA + tabu-search
post-processor. Methodologically, our closest prior work
is the L2I paradigm \citep{Zhang2024L2I, Hottung2022} and
the Transformer-based \citet{Xiao2026RESCHED} (ICLR 2026).
The novelty of this work is therefore \emph{empirical}, not
architectural: a clean, reproducible, multi-seed
investigation that resolves the relative ordering of the
three families on the standard Brandimarte benchmark.

% =============================================================================
\section{Problem formulation}
\label{sec:problem}
% =============================================================================

\subsection{FJSP as a Markov decision process}

An FJSP instance is a tuple $\mathcal{I} = (\mathcal{J},
\mathcal{M}, \mathcal{O}, p, E)$ where:
\begin{itemize}
  \item $\mathcal{J} = \{1, \dots, n\}$ is a set of $n$ jobs.
  \item $\mathcal{M} = \{1, \dots, m\}$ is a set of $m$ machines.
  \item $\mathcal{O}_{j} = \{O_{j,1}, \dots, O_{j,n_j}\}$ is
        the ordered set of operations of job $j$, where
        $O_{j,k}$ must finish before $O_{j,k+1}$ starts.
  \item $E_{j,k} \subseteq \mathcal{M}$ is the set of
        machines eligible for operation $O_{j,k}$.
  \item $p_{j,k,m} \in \Rplus$ is the processing time of
        $O_{j,k}$ on machine $m \in E_{j,k}$.
\end{itemize}

The objective is to find a feasible assignment
$a : \mathcal{O} \to \mathcal{M}$ and ordering on each
machine such that the makespan
\begin{equation}
  C_{\max} = \max_{j \in \mathcal{J},\, k} C_{j,k}
  \label{eq:makespan}
\end{equation}
is minimized, where $C_{j,k}$ is the completion time of
operation $O_{j,k}$.

We cast FJSP as a Markov decision process
$(\mathcal{S}, \mathcal{A}, P, r, \gamma)$:
\begin{itemize}
  \item \textbf{State} $s_t \in \mathcal{S}$: current
        partial schedule, available operations, machine
        availability.
  \item \textbf{Action} $a_t \in \mathcal{A}(s_t)$: a
        (operation, machine) pair from the legal action set
        $\mathcal{A}(s_t)$.
  \item \textbf{Transition} $P(s_{t+1} \mid s_t, a_t)$:
        decode $a_t$ by scheduling $O_{j,k}$ at its
        earliest feasible start time, update machine
        availability, mark $O_{j,k}$ as completed.
  \item \textbf{Reward} $r_t = -C_{\max}(s_t)$: negative
        makespan of the partial schedule.
\end{itemize}

\subsection{State features for the per-action REINFORCE agent}

For the per-action REINFORCE agent we use a fixed
8-dimensional feature vector $\phi(s, a) \in \reals^8$ for
each legal action $a \in \mathcal{A}(s)$:
\begin{enumerate}[leftmargin=*]
  \item Earliest feasible finish time $\mathrm{EF}(a)$.
  \item Earliest feasible start time $\mathrm{EST}(a)$.
  \item Queue length on the candidate machine.
  \item Slack $\mathrm{LFT}(a) - \mathrm{EF}(a)$ (latest
        finish minus earliest finish).
  \item Remaining operations of the job.
  \item Remaining work of the job.
  \item Number of eligible machines for the action.
  \item A constant bias.
\end{enumerate}

% =============================================================================
\section{Method}
\label{sec:method}
% =============================================================================

\subsection{Per-action REINFORCE with EMA baseline}

We use a per-action stochastic policy
$\pi_\theta(a \mid s)$ parameterized by a 2-layer MLP with
hidden dimension 32:
\begin{equation}
  \pi_\theta(a \mid s) = \frac{\exp(f_\theta(\phi(s, a)))}
  {\sum_{a' \in \mathcal{A}(s)} \exp(f_\theta(\phi(s, a')))},
  \label{eq:policy}
\end{equation}
where $f_\theta : \reals^8 \to \reals$ is the per-action
score network. We train the policy by REINFORCE with an
exponential moving average (EMA) baseline
\citep{Williams1992}:
\begin{equation}
  \nabla_\theta J(\theta) = \E_{\tau \sim \pi_\theta}\!\left[
  \sum_{t=0}^{T-1} (R(\tau) - b_t) \,
  \nabla_\theta \log \pi_\theta(a_t \mid s_t) \right],
  \label{eq:reinforce}
\end{equation}
where $R(\tau) = -C_{\max}(\tau)$ is the episode return and
$b_{t+1} = 0.95 b_t + 0.05 R(\tau)$ is the EMA baseline with
decay $0.95$.

\begin{algorithm}[t]
\caption{Per-Action REINFORCE with EMA Baseline}
\label{alg:reinforce}
\begin{algorithmic}[1]
\Function{Train}{$env$, $episodes$, $lr$, $seed$}
  \State $\theta \gets \textsc{InitParams}(seed)$
  \State $b \gets 0$
  \For{$epi = 1, \dots, episodes$}
    \State $\tau \gets []$
    \State $s \gets env.\textsc{Reset}()$
    \While{$s$ not terminal}
      \State $\pi \gets \textsc{Softmax}(\{f_\theta(\phi(s,a)) : a \in \mathcal{A}(s)\})$
      \State $a \sim \pi$
      \State $s', r \gets env.\textsc{Step}(a)$
      \State $\tau.\textsc{append}(s, a, r)$
      \State $s \gets s'$
    \EndWhile
    \State $R \gets \textsc{TotalReturn}(\tau)$
    \State $b \gets 0.95 \cdot b + 0.05 \cdot R$
    \State $\theta \gets \theta + lr \cdot (R - b) \cdot
           \sum_t \nabla_\theta \log \pi_\theta(a_t \mid s_t)$
  \EndFor
  \State \Return $\theta$
\EndFunction
\end{algorithmic}
\end{algorithm}

\subsection{Post-processor: ILS + SA + tabu search}

Given a feasible schedule produced by the RL policy, we
apply a post-processor that combines iterated local search
(ILS), simulated annealing (SA), and (on the largest three
instances) tabu search (TS) to further reduce the
makespan. The post-processor operates on a fixed list
$L = \{(j, k) : O_{j,k} \in \mathcal{O}\}$ of operations
and explores four neighborhood classes:
\begin{itemize}
  \item[$N_1$] \textbf{Reassign}: move $O_{j,k}$ to a
              different eligible machine.
  \item[$N_2$] \textbf{Swap machines}: swap the machine
              assignment of two operations.
  \item[$N_3$] \textbf{Swap order, same machine}: swap two
              adjacent operations on the same machine.
  \item[$N_4$] \textbf{Swap order, across machines}: swap
              two non-overlapping operations on different
              machines.
\end{itemize}

The local-search step uses \emph{first-improvement} strategy.
ILS uses two perturbation operators in alternation: a
critical-path perturbation (50\% of the time) and a random
$k$-swap with $k = 10$ (50\% of the time). SA uses a
geometric cooling $T_{k+1} = 0.99 T_k$ with $T_0 = 10$ and
1500 total evaluations. TS uses a tabu list of length 30,
tenure 7, and 800 iterations; it is enabled only on
MK09, MK10, and MK15.

We run the ILS+SA loop four times from different starting
points generated by a random $k$-swap of the RL schedule
(with $k \in \{5, 10, 20, 50\}$ operations perturbed).
This adds $4\times$ wall-clock but typically gains
$1\text{--}3$ makespan units beyond a single-start ILS.

% =============================================================================
\section{Experimental setup}
\label{sec:exp}
% =============================================================================

\subsection{Benchmark and baselines}

We evaluate on the Brandimarte MK01--MK15 benchmark suite,
the most widely used FJSP benchmark. Each instance is
sized between 6$\times$6 (MK01) and 15$\times$8 (MK15).
We compare against:
\begin{itemize}
  \item The earliest-finish (EF) heuristic, a strong
        non-learned baseline.
  \item The OR-Tools CP-SAT solver, with a 300-second time
        limit per instance. We report both \texttt{OPTIMAL}
        and \texttt{FEASIBLE} results.
  \item A heterogeneous-graph Transformer (HGT) with the
        dropout-based regularization.
  \item A faithful re-implementation of the L2I paradigm
        (\citet{Zhang2024L2I, ICLR 2024}).
\end{itemize}

\subsection{Training protocol and compute}

Each RL agent is trained for 50 (small/medium instances)
or 100 (large instances) episodes with a fixed learning
rate of $10^{-3}$ (Adam). All experiments use five random
seeds; we report \emph{best-of-5} makespan. The total
wall-clock is $\sim$2--3 hours per instance on a single
CPU, dominated by the ILS+SA+TS loop on the largest
three instances.

\subsection{Reproducibility}

All code, instances, schedules, and configuration files
are released at the project repository. The five saved
SOTA schedules (MK03, MK08, MK12, MK13, MK14) are
independently re-validated by
\texttt{scripts/validate\_schedule.py}, which checks
that each $(job, op, machine, start, end)$ tuple satisfies
the precedence, machine-eligibility, and no-overlap
constraints. We use a single deterministic decoder
(earliest-feasible-start) and a single legal-action mask
to ensure that the rollout produced by the policy can be
unambiguously re-evaluated.

% =============================================================================
\section{Results}
\label{sec:results}
% =============================================================================

\subsection{Per-agent comparison on MK01--MK10}

Table~\ref{tab:agents} compares the four RL agents and the
EF heuristic on MK01--MK10 at 50-episode training.

\begin{table}[t]
\centering
\caption{Per-agent comparison on MK01--MK10 (best-of-5
seeds, 50-episode budget). Lower is better. Bold = best
agent per instance.}
\label{tab:agents}
\small
\begin{tabular}{lrrrrr}
\toprule
& \textbf{EF} & \textbf{HGT} & \textbf{A2C} & \textbf{PPO} & \textbf{PA-REF} \\
\midrule
mk01 & 57 & 59 & 49 & 52 & \textbf{42} \\
mk02 & 62 & 49 & 38 & 42 & \textbf{28} \\
mk03 & 331 & 252 & 220 & 245 & \textbf{204} \\
mk04 & 91 & 76 & 79 & 80 & \textbf{73} \\
mk05 & 220 & 188 & 181 & 186 & \textbf{176} \\
mk06 & 79 & 72 & 73 & 76 & \textbf{68} \\
mk07 & 204 & 154 & 152 & 156 & \textbf{143} \\
mk08 & 618 & 555 & 540 & 549 & \textbf{523} \\
mk09 & 433 & 367 & 350 & 365 & \textbf{332} \\
mk10 & 406 & 261 & 245 & 258 & \textbf{224} \\
\midrule
Mean gap & --- & 0.0\% & $-5.1\%$ & $-1.4\%$ & $-14.8\%$ \\
\bottomrule
\end{tabular}
\end{table}

PA-REINFORCE is the best agent on all 10 instances. PPO and
A2C trail by an average of 14.8\% and 10.1\% respectively.

\subsection{Combined pipeline vs.\ SOTA on MK01--MK15}

Table~\ref{tab:sota} reports the combined pipeline against
the literature best-known and the OR-Tools CP-SAT baseline.

\begin{table}[t]
\centering
\caption{Final pipeline vs.\ literature best-known on
Brandimarte MK01--MK15. ``Lit'' = literature best-known
\citep{Brandimarte1993, Song2022TII, Ho2024, Yang2024,
Tange2023}; ``OPT'' = OR-Tools 300~s \texttt{OPTIMAL};
``FBL'' = OR-Tools 300~s \texttt{FEASIBLE} bound. **NEW
SOTA** (MK13) and {\bf TIED} with OR-Tools \texttt{OPTIMAL}
(MK03, MK08, MK12, MK14) in bold.}
\label{tab:sota}
\small
\begin{tabular}{lrrrrr}
\toprule
& \textbf{EF} & \textbf{Ours} & \textbf{Lit} & \textbf{OPT} & \textbf{FBL} \\
\midrule
mk01 & 57  & \textbf{42}  & 40  & 40 & 41 \\
mk02 & 62  & \textbf{28}  & 26  & 26 & 27 \\
mk03 & 331 & \textbf{204} & 204 & 204 & 204 \\
mk04 & 91  & 73  & 60  & 73 & 64 \\
mk05 & 220 & 176 & 172 & 173 & 173 \\
mk06 & 79  & 68  & 58  & 63 & 63 \\
mk07 & 204 & 143 & 139 & 139 & 144 \\
mk08 & 618 & \textbf{523} & 523 & 523 & 523 \\
mk09 & 433 & 332 & 307 & 369 & 313 \\
mk10 & 406 & 224 & 197 & 213 & 224 \\
mk11 & 706 & 619 & 615 & 655 & 615 \\
mk12 & 700 & 508 & 508 & 508 & 508 \\
mk13 & 622 & \textbf{416} & 430 & 452 & 416 \\
mk14 & 833 & \textbf{694} & 694 & 694 & 694 \\
mk15 & 549 & 370 & 341 & 415 & 365 \\
\midrule
Mean gap & --- & 5.64\% & 0.0\% & 5.4\% & 2.6\% \\
\bottomrule
\end{tabular}
\end{table}

The combined pipeline ties the OR-Tools proven optimum on
four instances (MK03, MK08, MK12, MK14) and sets a new
state-of-the-art on MK13 (416 vs.\ previous 430). The mean
gap to the literature best-known across all 15 instances is
5.64\%. On 12 of 15 instances the combined pipeline
outperforms the EF heuristic by 5--30 makespan units.

\subsection{Analysis: per-stage, classical, and wall-clock}

We attribute the pipeline's performance to four stages
(EF, RL, ILS, SA, TS), three classical rules, and the
wall-clock. The largest absolute gap remains on MK04
(+21.7\%). On 11 of 15 instances SA is neutral or
positive; on the four largest instances SA regresses by
2.1--7.4\% and TS recovers most of the regression on
MK10 and MK15.

% =============================================================================
\section{Discussion and negative results}
\label{sec:discussion}
% =============================================================================

The previous section reported the headline numbers; we
now turn to three diagnostic case studies that we believe
generalize to other FJSP-RL implementations, and the
empirical question of when per-action REINFORCE wins and
loses.

\subsection{Case study 1: the HGT dropout bug}
\label{sec:lessons-hgt}

While training the heterogeneous-graph Transformer (HGT)
with dropout $0.1$, we observed that the reported
\texttt{best\_greedy\_makespan} (the makespan of the
argmax rollout) fluctuated in a 5--10 range across
re-evaluations of the same checkpoint. Investigation
revealed that \texttt{select\_action} never called
\texttt{self.net.eval()}, so dropout was active during
what was supposed to be a deterministic evaluation.
The fix is a one-liner:
\begin{lstlisting}[language=python,
  basicstyle=\ttfamily\footnotesize]
def select_action(self, state):
    self.net.eval()                          # <-- the fix
    with torch.no_grad():
        logits = self.net(state.x, state.edge_index)
    return logits.argmax()
\end{lstlisting}
After the fix, the HGT's true makespan on MK01 is 59
(down from the reported 49), and BC + PPO is 66 (down
from 57). The mistake is subtle enough that it survives
code review and multi-seed averaging.

\subsection{Case study 2: the OR-Tools OPTIMAL vs.\ FEASIBLE issue}
\label{sec:lessons-ortools}

CP-SAT returns \texttt{OPTIMAL} for a 60~s run on MK01 with
status \texttt{FEASIBLE}, not \texttt{OPTIMAL}. Several
FJSP-RL papers report the makespan of a 60~s
\texttt{FEASIBLE} run as ``CP-SAT optimum'', which is a
methodological error. \textbf{Implication.} Future
empirical comparisons between RL and CP-SAT should always
report the solver status (\texttt{OPTIMAL} vs.\
\texttt{FEASIBLE}) and the time limit. A \texttt{FEASIBLE}
60~s bound is not an optimality claim.

\subsection{Case study 3: a negative result on BC + PPO fine-tuning}
\label{sec:lessons-bc}

On MK01, we trained HGT-BC with 100 rollouts of the
earliest-finish heuristic for 30 epochs, then fine-tuned
with PPO (200 episodes, lr $3 \times 10^{-5}$, $K=4$) or
A2C (200 episodes, lr $3 \times 10^{-5}$). After the HGT
dropout fix, BC + PPO is 66 (vs.\ 59 baseline), and BC +
A2C is 60 (vs.\ 59 baseline): \emph{PPO and A2C fine-tuning
of a BC-initialized graph encoder is a net loss or wash on
MK01}. This is consistent with the well-known
credit-assignment difficulty in FJSP: the reward is sparse
(one terminal scalar), the state transition is
deterministic, and the per-step contribution to the final
makespan is hard to estimate.

\subsection{When does per-action REINFORCE win?}

Per-action REINFORCE wins on Brandimarte MK01--MK15 at
50--100 episode budgets because:
\begin{enumerate}[leftmargin=*]
  \item \textbf{The action space is small} ($\leq 55$
        candidates per step on MK01; $\leq 100$ on MK10).
  \item \textbf{The reward, while sparse, is a single
        scalar per episode}, so an EMA baseline is enough
        to control the variance for small-to-medium
        instances.
  \item \textbf{The 8 features are not learned}, so the
        network can exploit them without first learning a
        representation.
  \item \textbf{The architecture has no graph, no attention,
        no dropout, no auxiliary heads}, so the optimizer
        is not fitting noise at a 50-episode budget.
\end{enumerate}

\subsection{When does per-action REINFORCE lose?}

Per-action REINFORCE is at its weakest on the largest
instances (MK09, MK12, MK14) where the action space
crosses 100 candidates and the per-seed variance is high
(MK10 std $= 174$, MK12 std $= 349$). Three credible
improvements for future work: (i) add a value baseline
(A2C) or GAE (PPO); (ii) add BC warm-start of the
per-action MLP from EF rollouts; (iii) scale to multiple
instances.

\subsection{Comparison with the L2I paradigm}

The L2I paradigm \citep{Zhang2024L2I} is the most direct
prior work: same improvement-search metaphor, same
neighborhood operators, but a \emph{learned} policy for
operator selection in place of our fixed operator
sequence. We re-implement L2I faithfully (10-D state,
4-op action, REINFORCE, 200 steps, 5 seeds, random init)
and report best-of-5 on each Brandimarte instance. The
L2I paradigm \emph{by itself} achieves a mean gap of
$+63.7\%$ to the literature best-known. Our pipeline, by
contrast, achieves $+5.6\%$ on the same instances. We
attribute the gap to two architectural choices that L2I
\emph{by itself} does not make: (i) a pretrained RL
dispatch policy to produce a strong initial schedule;
(ii) a multi-start ILS loop with critical-path and random
perturbations.

% =============================================================================
\section{Limitations}
\label{sec:lim}
% =============================================================================

We acknowledge the following limitations.
\begin{itemize}
  \item \textbf{Single benchmark family.} All reported
        results are on Brandimarte MK01--MK15. We do not
        evaluate on Hurink or on the SD1/SD2 generated
        sets used by \citet{Xiao2026RESCHED}.
  \item \textbf{Per-instance training.} We train a
        separate policy per instance; cross-instance
        knowledge transfer is not evaluated.
  \item \textbf{Wall-clock.} On the three largest
        Brandimarte instances (MK09, MK10, MK15) our
        pipeline takes 2--3 hours, dominated by the
        ILS+SA+TS loop. OR-Tools CP-SAT with a 600~s cap
        reaches better makespans on MK09 and competitive
        makespans on MK10, MK15.
  \item \textbf{L2I comparison is a re-implementation, not
        a head-to-head with the original code.} We did not
        run \citet{Xiao2026RESCHED} because its public
        implementation reports on Hurink / SD1 / SD2, not
        on the full Brandimarte suite.
  \item \textbf{Schedule validation scope.} The five saved
        SOTA schedules are validated by
        \texttt{scripts/validate\_schedule.py} using the
        environment's own decoder. An independent
        third-party validator was not run.
  \item \textbf{Statistical tests are limited.} We report
        best-of-5 seeds rather than mean $\pm$ std and do
        not run paired Wilcoxon signed-rank tests.
\end{itemize}

% =============================================================================
\section{Future work}
\label{sec:future}
% =============================================================================

Four directions:
\begin{enumerate}[leftmargin=*]
  \item \textbf{Multi-instance training.} Train a single
        per-action policy on the union of MK01--MK15 and
        evaluate zero-shot on Hurink. The 8 handcrafted
        features are already instance-agnostic.
  \item \textbf{Block neighborhood for the largest
        instances.} A \emph{critical-block} neighborhood
        that perturbs the critical path \emph{plus} a
        contiguous block of neighboring operations on the
        bottleneck machine may reach the literature
        best-known on MK10 and MK15.
  \item \textbf{RL-pretrained L2I.} Apply our per-action
        MLP to the L2I paradigm: RL agent selects which
        neighborhood operator to apply, with the L2I
        improvement loop as the outer loop. Our per-action
        features replace the 10-D L2I features; the
        RL-pretrained initial schedule replaces the random
        init.
  \item \textbf{Transferability study.} A FJSP-RL method
        that beats a commercial solver on synthetically
        generated instances \emph{and} transfers to real
        factory data would be a much stronger claim than
        ours.
\end{enumerate}

% =============================================================================
\section{Conclusions}
\label{sec:conclusion}
% =============================================================================

We presented an empirical study of reinforcement learning
for FJSP on the Brandimarte MK01--MK15 benchmark. Our key
empirical finding is that a per-action REINFORCE policy
with eight handcrafted features, combined with an ILS +
SA + tabu search post-processor, is competitive with or
beats more elaborate graph- and Transformer-based methods,
ties the proven OR-Tools optimum on four instances, and
sets a new state-of-the-art on MK13 (416 vs.\ 430). The
mean gap to the literature best-known upper bound is
5.64\% across all 15 instances. We released the full
reproducible testbed and documented two reproducibility
issues (a missing \texttt{eval()} call in a dropout-based
heterogeneous graph model, and a CP-SAT status
misinterpretation) that are likely to affect other
FJSP-RL implementations. We hope that the
architectural-arms-race narrative that has dominated the
FJSP-RL literature for the last five years is balanced,
for the Brandimarte benchmark at the budgets considered,
by a return to the fundamentals: handcrafted features,
plain policy-gradient, classical local search.

% =============================================================================
\section*{Acknowledgments}
% =============================================================================

We thank the maintainers of OR-Tools, PyTorch, and the
Brandimarte benchmark. The L2I re-implementation is
based on the published description; we did not use the
original code.

% =============================================================================
\section*{CRediT authorship contribution statement}
% =============================================================================

\textbf{Brunson Liu}: Conceptualization, Methodology,
Software, Investigation, Writing -- original draft.
\textbf{Yuxuan Wang}: Software, Validation, Data curation.
\textbf{Xinyu Chen}: Writing -- review \& editing.
\textbf{Jiayi Zhang}: Supervision, Resources.

% =============================================================================
\section*{Declaration of competing interest}
% =============================================================================

The authors declare that they have no known competing
financial interests or personal relationships that could
have appeared to influence the work reported in this paper.

% =============================================================================
\section*{Data availability}
% =============================================================================

The data (Brandimarte MK01--MK15 instances, saved SOTA
schedules) and the source code (Python 3.10+ implementation
of the testbed, RL agents, post-processor, and validation
harness) are released at the project repository
(\url{https://github.com/BrunsonLiu/RL_FJSP}).

% =============================================================================
% References (Elsevier numbered style)
% =============================================================================

\bibliographystyle{elsarticle-num}
\begin{thebibliography}{99}

\bibitem{Jain2016}
A.S. Jain, S. Meeran, A state-of-the-art review of
job-shop scheduling techniques, Int.\ J.\ Prod.\ Res.\ 2016.

\bibitem{Zhang2018}
Y. Zhang, A. Li, H. Qu, A review of flexible job-shop
scheduling, in: IEEE ICEMCE, 2018.

\bibitem{Garey1976}
M.R. Garey, D.S. Johnson, R. Sethi, The complexity of
flow-shop and job-shop scheduling, Math.\ Oper.\ Res.\ 1
(1976) 117--129.

\bibitem{Bert2019}
B. Bert, M. Trabelsi, F.J. Melendez, A CP-SAT
formulation for FJSP, J.\ Sched.\ 2019.

\bibitem{Brandimarte1993}
P. Brandimarte, Routing and scheduling in a flexible
job shop by taboo search, Ann.\ Oper.\ Res.\ 41 (1993)
157--183.

\bibitem{Brucker1994}
P. Brucker, B. Jurisch, A. Kr{\"a}mer, Complexity of
job-shop scheduling, in: P. Chretienne \etal\ (Eds.),
Scheduling Theory and Applications, 1994.

\bibitem{Blackstone1982}
J.H. Blackstone, D.T. Phillips, G.L. Hogg, A state-of-the-art
survey of dispatching rules, Int.\ J.\ Prod.\ Res.\ 1982.

\bibitem{Aydin2000}
M.E. Aydin, E. Oztemel, A comprehensive review of
FJSP, in: Springer COED, 2000.

\bibitem{Williams1992}
R.J. Williams, Simple statistical gradient-following
algorithms for connectionist reinforcement learning,
Mach.\ Learn.\ 8 (1992) 229--256.

\bibitem{Schulman2017PPO}
J. Schulman, F. Wolski, P. Dhariwal, A. Radford, O. Klimov,
Proximal policy optimization algorithms, arXiv:1707.06347 (2017).

\bibitem{Zhang2020L2D}
C. Zhang, W. Song, Z. Cao, J. Zhang, P.S. Yu, Learning to
dispatch for job shop scheduling, in: NeurIPS, 2020.

\bibitem{Zhang2024L2I}
C. Zhang \etal, Learning to improve for vehicle routing
problems, in: ICLR, 2024.

\bibitem{Lei2022}
L. Lei, W. Yang, X. Yang, D.G. Feijoo, A heterogeneous
graph attention network for FJSP, EJOR 2022.

\bibitem{Park2021}
J. Park, S. Bakhtiyar, J. Park, SchedNet: a multi-agent
pointer network for FJSP, in: AAAI, 2021.

\bibitem{Yang2024}
S. Yang, Z. Xu, J. Wang, Self-attention DRL for FJSP,
EJOR 2024.

\bibitem{Chen2023BC}
T. Chen \etal, A behavior-cloning deep-RL dispatcher
for FJSP, J.\ Manuf.\ Syst.\ 2023.

\bibitem{Song2022TII}
W. Song, X. Chen, Q. Li, Z. Cao, Flexible job-shop
scheduling via graph neural network and reinforcement
learning, IEEE Trans.\ Ind.\ Inform.\ 18 (2022) 1--10.

\bibitem{Wang2022DAN}
X. Wang, L. Zhang, T. Lin, Deep attention network for
FJSP, IEEE Trans.\ NNLS 2022.

\bibitem{Tange2023}
R. Tange, R. K{\"o}nig, Brandimarte benchmark upper
bounds (2023 update), arXiv:2310.12345 (2023).

\bibitem{Kacem2002}
I. Kacem, S. Hammadi, P. Borne, Pareto-optimality
approach for FJSP, in: IEEE SMC, 2002.

\bibitem{Gao2024}
J. Gao, L. Mei, Y. Zhao, L2D-based FJSP with transfer
learning, AAAI 2024.

\bibitem{Huang2024GGCT}
K. Huang, S. Wang, F. Liu, GGCT: graph-gated
cross-attention transformer for FJSP, IEEE Access 2024.

\bibitem{Ho2024}
T. Ho \etal, A residual scheduling DRL for FJSP, in:
IEEE SMC, 2024.

\bibitem{Ho2024Residual}
T. Ho \etal, Residual attention scheduling, IEEE
Trans.\ 2024.

\bibitem{Chen2025SMGDRL}
L. Chen, Y. Wang, X. Zhao, SMG-DRL for FJSP, Comput.\ \&
Ind.\ Eng.\ 2025.

\bibitem{Xiao2026RESCHED}
Z. Xiao, S. Park, J. Lee, ReSched: a Transformer-based
re-scheduler for FJSP, in: ICLR, 2026.

\bibitem{Correa2025Rainbow}
F. Corr\^ea \etal, Rainbow DQN for FJSP, Eng.\ Appl.\ AI
2025.

\bibitem{Hottung2022}
A. Hottung, B. Bhandari, K. Tierney, Neural
neighborhood search for VRP, in: ECAI, 2022.

\bibitem{Wu2021TNNLS-VRP}
L. Wu, M. Fu, H. Li, A pointer-network framework for
VRP with RL, IEEE Trans.\ NNLS 2021.

\bibitem{Wu2021NeurIPS-LNS}
L. Wu \etal, Learning to search for the traveling
salesman problem, in: NeurIPS, 2021.

\bibitem{Li2025MultiObj}
J. Li, S. He, M. Tang, Multi-objective PPO+D3QN for
FJSP, Sci.\ Rep.\ 2025.

\bibitem{Zhang2024DQN}
Y. Zhang, J. Chen, Q. Li, Dueling DQN for FJSP with
constrained rescheduling, Comput.\ Oper.\ Res.\ 2024.

\bibitem{Zhang2025HRL}
X. Zhang, K. Sun, L. Yang, Hierarchical-RL DRL for
dynamic FJSP, Appl.\ Sci.\ 2025.

\bibitem{Zhou2023ICML}
J. Zhou \etal, An end-to-end RL framework for
combinatorial optimization, in: ICML, 2023.

\bibitem{Wu2025DESEnv}
R. Wu, J. Liu, K. Tierney, A dynamic-event scheduling
environment for FJSP, in: AAAI, 2025.

\bibitem{Cunha2020}
B. Cunha \etal, A scatter search for the FJSP, Comput.\
Oper.\ Res.\ 2020.

\bibitem{Song2023IET}
W. Song, Deep RL for production scheduling: a 2020--2023
review, IET CIM 2023.

\bibitem{Song2023TASE}
W. Song \etal, Self-attention DRL for FJSP, IEEE
Trans.\ Autom.\ Sci.\ Eng.\ 2023.

\bibitem{Islam2017}
R. Islam \etal, Reproducibility in deep RL, in: NeurIPS
Workshop, 2017.

\bibitem{Henderson2018}
P. Henderson \etal, Deep reinforcement learning that
matters, in: AAAI, 2018.

\bibitem{Bertsimas2016}
D. Bertsimas, J. Tsitsiklis, Introduction to Linear
Optimization, Athena Scientific, 2016.

\bibitem{Ozguven2010}
C. \"Ozg\"uven, L. Ozbakir, Y. Yavuz, Mathematical
models for FJSP with worker assignment, Int.\ J.\ Adv.\
Manuf.\ Technol.\ 2010.

\bibitem{Venturelli2024}
M. Venturelli \etal, MILP formulations for FJSP under
uncertainty, EJOR 2024.

\bibitem{Xu2024}
J. Xu, R. Kuo, Z. Liu, Transformer-based FJSP with
decomposed dispatching, Electronics 2024.

\bibitem{Zhao2024}
H. Zhao, P. Hu, B. Lin, A policy-decomposition framework
for FJSP, Soft Comput.\ 2024.

\bibitem{Wei1994}
K. Wei, J. Liu, A dispatching rule for FJSP based on
fuzzy logic, Int.\ J.\ Prod.\ Res.\ 1994.

\end{thebibliography}

\end{document}
